\section{Background}
\subsection{Managed and Native Allocation}
ART garbage collection reasons about managed objects.  Native allocations may be
reported to the runtime through Android allocation-notification mechanisms, but
such notifications are not equivalent to ownership of native objects or complete
knowledge of allocator caches, mappings, sharing, and page residency.  Native
allocation rate is likewise not RSS growth: lifetime, frees, reuse, page
commitment, file mappings, and reclaim all intervene.  HMS consequently treats
managed allocation, sampled native allocation, and RSS as distinct observations.

\subsection{Linux and Android Memory Pressure}
Linux memory cgroups account and limit groups of processes and expose a
\texttt{memory.reclaim} interface for proactive, cgroup-scoped reclaim.  PSI means
\emph{Pressure Stall Information} and quantifies time in which tasks stall for
resources.  Android's LMKD consumes pressure and process-importance information;
zRAM provides compressed swap; MGLRU and DAMON/DAMOS provide additional
working-set and proactive-reclaim mechanisms.  HMS is a proposed trigger and
budgeting layer above these facilities, not a replacement for them.

\subsection{Memory Tagging}
Arm Memory Tagging Extension (MTE) detects classes of spatial and temporal memory
safety errors.  An MTE fault is not a normal heap-pressure signal, and this paper
does not use it as a predictor.  Production mode, overhead, causality, and fallback
would require a separate device study.  MTE is therefore outside the evaluated HMS
control path.
